% Báo cáo thực nghiệm; số liệu lấy từ results/validation_*.json.
\documentclass[11pt,a4paper]{article}
\usepackage[left=2cm,right=2cm,top=2.1cm,bottom=2cm,headheight=15pt]{geometry}
\usepackage{iftex}
\ifPDFTeX
  \usepackage[T5]{fontenc}
  \usepackage[utf8]{inputenc}
  \usepackage[vietnamese]{babel}
\else
  \usepackage{fontspec}
  \setmainfont{TeX Gyre Termes}
  \setsansfont{TeX Gyre Heros}
  \setmonofont{Latin Modern Mono}
\fi
\usepackage{amsmath,array,tabularx,booktabs}
\usepackage[table]{xcolor}
\usepackage{graphicx,enumitem,fancyhdr,listings,tikz,titlesec,caption,microtype}
\usetikzlibrary{arrows.meta,positioning}
\usepackage[unicode,hidelinks]{hyperref}
\definecolor{accent}{HTML}{13558C}
\definecolor{muted}{HTML}{51606C}
\definecolor{light}{HTML}{F2F6FA}
\definecolor{line}{HTML}{CCD9E4}
\setlength{\parindent}{0pt}
\setlength{\parskip}{5.5pt}
\linespread{1.05}
\renewcommand{\arraystretch}{1.18}
\setlist[itemize]{leftmargin=1.4em,itemsep=2pt,topsep=3pt}
\newcolumntype{Y}{>{\raggedright\arraybackslash}X}
\titleformat{\section}{\Large\sffamily\bfseries\color{accent}}{\thesection}{0.7em}{}[\vspace{-2mm}\color{line}\titlerule]
\titlespacing*{\section}{0pt}{13pt}{8pt}
\captionsetup{font=small,labelfont={bf,color=accent},skip=5pt}
\pagestyle{fancy}\fancyhf{}
\fancyhead[L]{\small\sffamily\color{muted}Tương tác người -- robot\quad/\quad Bài thực hành 02}
\fancyhead[R]{\small\sffamily\color{muted}\MSSV}
\fancyfoot[L]{\small\sffamily\color{muted}Nguyễn Hoàng Thiện}
\fancyfoot[R]{\small\sffamily\color{muted}Trang \thepage}
\renewcommand{\headrulewidth}{0.4pt}
\renewcommand{\headrule}{\hbox to\headwidth{\color{line}\leaders\hrule height \headrulewidth\hfill}}
\lstset{basicstyle=\ttfamily\small,backgroundcolor=\color{light},frame=single,
rulecolor=\color{line},breaklines=true,columns=fullflexible,keepspaces=true,
showstringspaces=false,aboveskip=7pt,belowskip=7pt,framesep=6pt}
\newcommand{\HoTen}{Nguyễn Hoàng Thiện}
\newcommand{\MSSV}{23020774}
\newcommand{\Anh}[1]{\IfFileExists{#1}{\includegraphics[width=\linewidth]{#1}}{\fbox{\parbox[c][3cm][c]{0.92\linewidth}{\centering Ảnh thực nghiệm đi kèm trong thư mục \texttt{results}.\\Trình biên dịch cần truy cập tệp ảnh.}}}}
\begin{document}
\begin{titlepage}
\thispagestyle{empty}
\centering
\vspace*{5mm}
{\large\sffamily\bfseries ĐẠI HỌC QUỐC GIA HÀ NỘI\par}
\vspace{2mm}
{\large\sffamily\bfseries TRƯỜNG ĐẠI HỌC CÔNG NGHỆ\par}
\vspace{5mm}{\color{accent}\rule{0.62\linewidth}{0.8pt}}\par
\vspace{12mm}
\IfFileExists{results/cover_logo_original.png}{\includegraphics[width=3.4cm]{results/cover_logo_original.png}}{\rule{3.4cm}{0pt}}\par
\vspace{16mm}
{\Large\sffamily\bfseries\color{accent} BÁO CÁO THỰC HÀNH 02\par}
\vspace{4mm}
{\large\sffamily Môn học: Tương tác người -- robot\par}
\vspace{13mm}
{\LARGE\sffamily\bfseries Điều khiển UR3e bằng LLM\par}
\vspace{3mm}
{\Large\sffamily\bfseries và Skill-based Planning\par}
\vspace{6mm}{\color{accent}\rule{0.18\linewidth}{1pt}}\par
\vfill
\begin{tabular}{@{}ll@{}}
\textcolor{accent}{\textbf{Sinh viên thực hiện:}} & \HoTen\\[3mm]
\textcolor{accent}{\textbf{Mã số sinh viên:}} & \MSSV\\[3mm]
\textcolor{accent}{\textbf{Giảng viên hướng dẫn:}} & Nguyễn Quốc Bảo
\end{tabular}
\vfill
{\color{muted}Hà Nội, tháng 9 năm 2026\par}
\end{titlepage}
\setcounter{page}{1}

\noindent\colorbox{light}{\parbox{\dimexpr\linewidth-2\fboxsep\relax}{\small
\textbf{\textcolor{accent}{Tóm tắt.}} Trong bài thực hành này, em xây dựng
hệ thống điều khiển UR3e mô phỏng bằng câu lệnh tiếng Việt và tiếng Anh.
LLM đề xuất thứ tự các skill, còn chương trình kiểm tra kế hoạch và dùng
MoveIt~2 để lập quỹ đạo. Với MSSV \MSSV, nhiệm vụ cá nhân là đưa khối vàng
vào A, khối đỏ vào B và khối xanh dương vào C. Lần chạy theo MSSV hoàn thành
7/7 skill; video và log được lưu cùng mã nguồn.}}
\vspace{4mm}

\section{Mục tiêu và môi trường triển khai}
Mục tiêu của bài là để robot hiểu một yêu cầu như ``đưa khối đỏ vào vùng B''
và tự thực hiện đúng chuỗi gắp, đặt, về vị trí chờ. Em triển khai chương trình
trên Ubuntu 22.04.5 trong WSL2, với ROS~2 Humble, Gazebo Classic 11.10.2,
MoveIt~2 phiên bản 2.5.10 và mô hình UR3e từ \texttt{ur\_description} 2.13.0.
Phần LLM được kết nối qua 9Router 0.5.91; các lần kiểm thử dùng model
\texttt{oc/mimo-v2.5-free}.

Em giới hạn đầu ra của LLM ở mức kế hoạch tác vụ: chọn \texttt{pick},
\texttt{place} hoặc \texttt{home} cùng tên vật và vùng đích. Chương trình mới
là nơi kiểm tra kế hoạch và yêu cầu MoveIt~2 tính đường đi. Cách chia này giúp
robot xử lý nhiều cách diễn đạt tiếng Việt, tiếng Anh mà không cho LLM gửi
lệnh khớp trực tiếp.

Trong Gazebo, em bố trí một bàn thao tác, ba khối cạnh 40~mm và ba vùng đặt
vuông cạnh 80~mm quanh UR3e. Tọa độ ban đầu và tâm vùng được khai báo trong
\texttt{scene.yaml}; khi chạy, chương trình đọc lại vị trí vật từ simulator.
Camera toàn cảnh chỉ ghi ảnh và video làm bằng chứng, không tham gia nhận dạng
hay xác định vị trí vật.

\section{Cấu trúc chương trình, node và package chính}
Em đặt mã nguồn trong package \texttt{ur3\_llm\_control}. Điểm vào là node
\texttt{skill\_executor}: node nhận câu lệnh ở CLI, gọi planner, kiểm tra kế
hoạch rồi chạy các skill theo thứ tự. Planner và validator là hai mô-đun Python
được node này sử dụng, không phải hai node ROS chạy độc lập. Những thành phần
chính và vai trò của chúng được tóm tắt ở bảng dưới.

\begin{tabularx}{\linewidth}{@{}p{4.25cm}Y@{}}
\toprule\textbf{Thành phần} & \textbf{Vai trò trong chương trình}\\\midrule
\texttt{llm\_planner.py} & Gửi prompt, trạng thái vùng đặt và ánh xạ MSSV tới 9Router; nhận JSON và lưu thông tin model/token.\\
\texttt{task\_validator.py} & Kiểm tra schema, allowlist, tham số và thứ tự gắp--thả; phát hiện vùng bị chiếm.\\
\texttt{skill\_executor.py} & Điều phối tuần tự, dừng khi lỗi; ghi log, tọa độ và ảnh sau từng skill.\\
\texttt{robot\_skills.py} & Gọi dịch vụ/action MoveIt, cập nhật planning scene và xác nhận kết quả từ Gazebo.\\
\texttt{skill\_grasp.cpp} & Plugin Gazebo tạo/tháo khớp cố định giữa đầu hút và vật ở gần.\\
\texttt{move\_group} & Bộ giải KDL, lập kế hoạch OMPL/RRTConnect và đường đi Cartesian.\\
\texttt{gazebo\_ros2\_control} & Kết nối robot mô phỏng với \texttt{joint\_trajectory\_controller}; phát trạng thái khớp.\\
\texttt{robot\_state\_publisher} & Phát cây TF từ mô hình robot và trạng thái khớp.\\
\bottomrule
\end{tabularx}

File \texttt{llm\_robot.launch.py} khởi động Gazebo, robot, controller và
MoveIt. Em đặt cấu hình trong \texttt{config}, prompt trong \texttt{prompts},
mô hình đầu hút trong \texttt{urdf} và cảnh trong \texttt{worlds}. Các phiên
chạy dùng \texttt{ROS\_DOMAIN\_ID=42} để tách khỏi những ROS graph khác trên máy.

\newpage
\section{Các robot skill đã xây dựng}
\begin{tabularx}{\linewidth}{@{}>{\raggedright\arraybackslash}p{3.5cm}Y@{}}
\toprule\textbf{Skill} & \textbf{Các bước và điều kiện thực thi}\\\midrule
\texttt{home()} & Kiểm tra tay gắp rỗng; dùng MoveIt đưa sáu khớp về tư thế chờ trong cấu hình.\\
\texttt{pick(object)} & Kiểm tra vật hợp lệ và tay gắp rỗng; đọc vị trí thực, tiếp cận từ trên cao, hạ đầu hút, gắn vật rồi nâng 100~mm. Xác nhận độ cao vật tăng ít nhất 60~mm.\\
\texttt{place(object, zone)} & Kiểm tra đang giữ đúng vật và vùng đích trống; di chuyển đến trên vùng, hạ xuống, tháo khớp gắp và rút đầu hút. Kiểm tra sai số vị trí bằng trạng thái Gazebo.\\
\bottomrule
\end{tabularx}

Em tách mỗi thao tác gắp--đặt thành hai kiểu chuyển động. Đoạn đi tới vị trí
tiếp cận được MoveIt lập kế hoạch bằng KDL và OMPL/RRTConnect qua
\texttt{/move\_action}. Khi đầu hút ở gần vật hoặc vùng đặt, em dùng đường
Cartesian để hạ, nâng và rút theo từng bước 5~mm. Chương trình chỉ gửi quỹ
đạo này đi khi MoveIt tính được đủ 100\% đường đi, đã tham số hóa thời gian và
kiểm tra va chạm lại trên các mẫu quỹ đạo. Hệ số vận tốc/gia tốc lần lượt là
0,18/0,15 cho đường OMPL và 0,15/0,10 cho đường Cartesian; quỹ đạo sau cùng
đi qua action \texttt{/execute\_trajectory}.

Trong Gazebo, đầu hút giữ vật bằng một khớp cố định và chỉ gắn khi đầu hút cách
tâm vật không quá 35~mm. Cùng lúc đó, em đánh dấu vật là \emph{attached
collision object} trong planning scene của MoveIt để bộ lập kế hoạch biết vật
đang đi cùng robot. Lúc bắt đầu nhấc, vật vẫn tiếp xúc với bàn nên em tạm cho
phép đúng cặp va chạm vật--bàn; sau khi nâng lên, kiểm tra cặp này được bật
lại. Robot vẫn phải tránh bàn, các vật khác và tự va chạm trong suốt thao tác.

Để biết một tác vụ thất bại ở đâu, em ghi riêng các mã lỗi như
\texttt{INVALID\_OBJECT}, \texttt{INVALID\_ZONE}, \texttt{PLANNING\_FAILED},
\texttt{GRASP\_FAILED} và \texttt{PRECONDITION\_FAILED}. Log của từng skill
gồm \texttt{SUCCESS} hoặc \texttt{FAILED} kèm nguyên nhân; nếu một skill lỗi,
executor không chạy tiếp các bước sau.

\newpage
\section{Luồng xử lý từ câu lệnh đến robot}
\begin{figure}[htbp]
\centering
\begin{tikzpicture}[
  stage/.style={draw=black,fill=white,align=center,text width=7.0cm,
    minimum height=8mm,font=\small,inner sep=3pt},
  reject/.style={draw=black,dashed,fill=white,align=center,text width=3.1cm,
    minimum height=8mm,font=\small,inner sep=3pt},
  flow/.style={-{Latex[length=2.2mm]},line width=0.55pt},
  node distance=4.5mm]
\node[stage] (user) {\textbf{1. Câu lệnh tự nhiên}\quad CLI nhận nguyên văn tiếng Việt/Anh};
\node[stage,below=of user] (llm) {\textbf{2. 9Router / LLM}\quad Prompt + vị trí vật + ánh xạ MSSV};
\node[stage,below=of llm] (json) {\textbf{3. Kế hoạch JSON}\quad Danh sách \texttt{pick}, \texttt{place}, \texttt{home}};
\node[stage,below=of json] (valid) {\textbf{4. Validator}\quad Schema, allowlist, thứ tự và vùng trống};
\node[stage,below=of valid] (exec) {\textbf{5. Skill Executor}\quad Chạy tuần tự, kiểm tra tiền điều kiện};
\node[stage,below=of exec] (moveit) {\textbf{6. MoveIt~2}\quad IK, đường đi, va chạm, giới hạn khớp};
\node[stage,below=of moveit] (robot) {\textbf{7. Controller + Gazebo/UR3e}\quad Thực thi quỹ đạo, đọc lại trạng thái};
\node[reject,right=7mm of valid] (badplan) {Sai kế hoạch: từ chối\\không chuyển động};
\node[reject,right=7mm of moveit] (badmotion) {Lỗi lập kế hoạch/skill:\\dừng và ghi log};
\draw[flow] (user)--(llm);
\draw[flow] (llm)--(json);
\draw[flow] (json)--(valid);
\draw[flow] (valid)--(exec);
\draw[flow] (exec)--(moveit);
\draw[flow] (moveit)--(robot);
\draw[flow,dashed] (valid)--(badplan);
\draw[flow,dashed] (moveit)--(badmotion);
\end{tikzpicture}
\caption{Luồng xử lý một yêu cầu và vị trí chương trình dừng khi kế hoạch hoặc chuyển động có lỗi.}
\label{fig:flow}
\end{figure}

\textbf{Nhận yêu cầu.} Khi em nhập một câu lệnh bằng tùy chọn
\texttt{--command}, node \texttt{skill\_executor} đọc vị trí ba vật từ
\texttt{/gazebo/get\_entity\_state} và đối chiếu với các vùng trong
\texttt{scene.yaml}. Từ đó, chương trình biết vùng nào đang trống trước khi
lập kế hoạch. Một ảnh toàn cảnh cũng được chụp ở thời điểm này để so sánh với
kết quả, nhưng ảnh không được dùng làm dữ liệu điều khiển.

\textbf{Hiểu lệnh và lập kế hoạch.} \texttt{llm\_planner.py} chuyển nguyên văn
câu lệnh, vị trí hiện tại và mô tả ba skill cho model qua 9Router. Nếu lệnh
nhắc đến MSSV, planner gửi thêm ánh xạ đích đã tính từ hai chữ số cuối. LLM
chỉ trả lời \emph{làm gì, với vật nào và đặt ở đâu}; nó không tạo tọa độ khớp
hay vận tốc. Chẳng hạn, với câu lệnh của T1, phản hồi có dạng:
\begin{lstlisting}
{"plan": [
  {"skill": "pick", "object": "red_cube"},
  {"skill": "place", "object": "red_cube", "zone": "zone_b"},
  {"skill": "home"}
]}
\end{lstlisting}

\textbf{Duyệt kế hoạch.} JSON được kiểm tra cú pháp và khóa trùng trước khi
validator xét nội dung. Kế hoạch chỉ được có tối đa 25 bước; mỗi skill phải
có đúng các trường cần thiết, tên vật và vùng phải thuộc danh sách cho phép.
Validator còn theo dõi trạng thái tay gắp: không thể \texttt{place} trước
\texttt{pick}, cầm hai vật cùng lúc hoặc đặt vào vùng đã có vật. Bước cuối
phải là \texttt{home} khi tay gắp đã rỗng. Nếu không đạt, chương trình báo
\texttt{FAILED} ngay tại đây; T4 là một ví dụ robot không hề chuyển động.

\textbf{Chạy từng thao tác.} Sau khi kế hoạch hợp lệ, executor thực hiện theo
đúng thứ tự và kiểm tra lại tiền điều kiện trước mỗi skill. Với
\texttt{pick}, robot tiếp cận vị trí vật vừa đọc từ Gazebo, hạ đầu hút, gắn
vật rồi nâng lên. Với \texttt{place}, robot đi tới vùng đích, hạ vật, tháo
khớp gắp và rút đầu hút. \texttt{home} chỉ chạy khi tay gắp rỗng.

\textbf{Lập quỹ đạo và xác nhận kết quả.} Mỗi chuyển động của skill đi qua
MoveIt~2 để giải động học nghịch, lập đường, kiểm tra va chạm và giới hạn
khớp. Quỹ đạo được gửi tới \texttt{joint\_trajectory\_controller}; Gazebo
cập nhật tư thế robot và vật. Sau mỗi skill, em đọc lại trạng thái từ Gazebo
để kiểm tra thao tác thực sự đạt điều kiện đích rồi mới ghi \texttt{SUCCESS}.
Nếu MoveIt hoặc skill báo lỗi, executor lưu nguyên nhân, trạng thái và ảnh
tại thời điểm lỗi, đồng thời dừng phần kế hoạch còn lại.

\newpage
\section{Trích đoạn mã giải thích luồng xử lý}
Các hình dưới đây được dựng từ \emph{mã nguồn thật} trong package, giữ nguyên
nội dung và số dòng; dấu ba chấm chỉ những dòng không liên quan được lược bớt.
Các dòng nền vàng là chỗ quyết định dữ liệu đi qua từng tầng, giúp đối chiếu
sơ đồ Hình~\ref{fig:flow} với phần triển khai.

\begin{figure}[htbp]\centering
\Anh{results/code_llm_planner.png}
\caption{\texttt{llm\_planner.py}: đọc MSSV, tạo ánh xạ trong ngữ cảnh LLM và phân tích phản hồi JSON.}
\label{fig:code-planner}
\end{figure}
\noindent Hình~\ref{fig:code-planner} cho thấy câu lệnh không được đổi bằng
những câu mẫu cố định: planner gửi ngữ cảnh \texttt{student\_assignment} và
\texttt{locations} tới model, rồi chỉ nhận lại JSON để chuyển sang validator.

\begin{figure}[htbp]\centering
\Anh{results/code_task_validator.png}
\caption{\texttt{task\_validator.py}: allowlist tham số của skill và kiểm tra vùng đích đã bị chiếm.}
\label{fig:code-validator}
\end{figure}
\noindent Hình~\ref{fig:code-validator} minh họa ranh giới an toàn: mỗi skill
chỉ nhận đúng bộ tham số cho phép; vùng bị chiếm bị từ chối trước khi robot
nhận quỹ đạo.

\begin{figure}[htbp]\centering
\Anh{results/code_skill_executor.png}
\caption{\texttt{skill\_executor.py}: gọi các robot skill tuần tự, ghi trạng thái và dừng khi có lỗi.}
\label{fig:code-executor}
\end{figure}
\noindent Hình~\ref{fig:code-executor} cho thấy executor gọi \texttt{home},
\texttt{pick} hoặc \texttt{place} theo từng bước đã duyệt; nó chụp ảnh và đọc
lại tọa độ sau mỗi skill để tạo bằng chứng thực nghiệm.

\clearpage
\section{Kết quả chạy thử}
Bộ kiểm thử xác nhận được chạy ngày 28/09/2026 sau khi hiệu chỉnh hệ thống.
Các ca hợp lệ dùng LLM thật qua 9Router, không nạp sẵn JSON thử nghiệm. T1--T3
thực thi tuần tự trên cùng cảnh; T6 và T7 mỗi ca khởi động từ cảnh ban đầu.

\begingroup\small
\begin{tabularx}{\linewidth}{@{}p{0.6cm}Yp{3.55cm}r@{}}
\toprule\textbf{Ca} & \textbf{Câu lệnh / tình huống} & \textbf{Kết quả quan sát} & \textbf{Giây}\\\midrule
T1 & Đưa khối màu đỏ vào vùng B. & 3/3 skill thành công; đỏ ở B. & 37,6\\
T2 & Hãy lấy khối màu vàng và đặt nó vào ô A. & 3/3 skill thành công; vàng ở A. & 41,5\\
T3 & Move the blue cube to zone C. & 3/3 skill thành công; xanh ở C. & 36,7\\
T4 & Đưa khối màu tím vào vùng D. & LLM trả kế hoạch rỗng; validator từ chối; không chuyển động. & 8,2\\
T5 & Arrange all objects according to my student ID. (Trước khi khai báo MSSV.) & Từ chối kế hoạch rỗng; không tự tạo MSSV hoặc nhiệm vụ. & 5,8\\
T6 & Đặt khối đỏ vào vùng A, khối vàng vào vùng B và khối xanh dương vào vùng C, rồi về home. & 7/7 skill thành công; ba vật đúng đích trong một tác vụ. & 87,4\\
T7 & Arrange all objects according to my student ID. (MSSV 23020774.) & 7/7 skill thành công; A--vàng, B--đỏ, C--xanh. & 84,6\\
\bottomrule
\end{tabularx}
\endgroup

Thời gian là thời gian thực trên máy, gồm gọi LLM, thực thi và ghi bằng chứng;
không phải thời gian riêng của thuật toán lập kế hoạch. Năm tác vụ hợp lệ hoàn
thành 23/23 skill. Cả 59 chuyển động con trả mã MoveIt \texttt{SUCCESS}; 36 đoạn
Cartesian đều đạt \texttt{fraction = 1.0}. Hai tác vụ không hợp lệ không gọi
skill và không làm thay đổi tọa độ các vật. T6 kiểm tra một lệnh phối hợp ba
vật với đích nêu trực tiếp; T7 mới là demo cá nhân hóa theo MSSV thật.

\textbf{Đối chiếu vị trí đích.} Hệ tọa độ \texttt{world}, đơn vị mét. Tâm vật
khi bắt đầu là $(0{,}26;-0{,}20;0{,}12)$, $(0{,}26;0;0{,}12)$ và
$(0{,}26;0{,}20;0{,}12)$ tương ứng đỏ, vàng, xanh dương.
\begin{center}\small
\begin{tabular}{@{}lllr@{}}
\toprule\textbf{Vật / đích} & \textbf{Tọa độ yêu cầu} & \textbf{Tọa độ đo sau tác vụ} & \textbf{Sai số (mm)}\\\midrule
Đỏ / B & $(0{,}40;0;0{,}12)$ & $(0{,}39998;-0{,}00003;0{,}12000)$ & 0,031\\
Vàng / A & $(0{,}40;-0{,}20;0{,}12)$ & $(0{,}39997;-0{,}20003;0{,}12000)$ & 0,040\\
Xanh / C & $(0{,}40;0{,}20;0{,}12)$ & $(0{,}40001;0{,}19996;0{,}12000)$ & 0,037\\
\bottomrule
\end{tabular}
\end{center}
Sai số tính bằng khoảng cách Euclid giữa tâm vật đo được và tâm đích. Đây là
kết quả của mô phỏng đầu hút lý tưởng, không thể suy ra độ chính xác robot thật.

\textbf{Kiểm thử validator.} 19 bài kiểm thử tự động đều đạt, bao gồm chín cặp
vật--vùng hợp lệ, schema/JSON sai, lệnh khớp chèn thêm, thứ tự thao tác sai,
vùng bị chiếm và đủ sáu ánh xạ MSSV. Log: \texttt{validator\_tests.log}
trong thư mục \texttt{results}.

\textbf{Các lỗi đã xử lý trong quá trình triển khai.} Đã sửa việc nạp YAML động
học để MoveIt nhận KDL, gắn đầu hút vào link có quán tính để Gazebo không loại
bỏ khi chuyển URDF, và tránh xóa vật hai lần khi chuyển sang trạng thái attached.
Đã giảm sai lệch cho phép khi tham số hóa quỹ đạo, chọn góc IK tương đương gần
trạng thái hiện tại và xử lý tiếp xúc vật--bàn ở bước nâng. Các lần chạy lỗi
trước hiệu chỉnh vẫn được lưu; bảng trên chỉ thống kê bộ xác nhận cuối, không
khẳng định tỷ lệ thành công tuyệt đối trong mọi điều kiện.
Một lần thử T7 dừng trước chuyển động do 9Router trả HTTP 502; sau khi khôi phục
kết nối, T7 được chạy lại từ cảnh ban đầu và hoàn thành. Log lỗi cũng được giữ
trong thư mục \texttt{results}.

\newpage
\section{Hình ảnh thực nghiệm}
\begin{center}
\begin{minipage}[t]{0.485\linewidth}\centering
\Anh{results/validation_T1.before_color_crop.png}
\small\textbf{(a)} Trước T1: ba khối ở vị trí ban đầu.
\end{minipage}\hfill
\begin{minipage}[t]{0.485\linewidth}\centering
\Anh{results/validation_T1.step1_color_crop.png}
\small\textbf{(b)} T1: khối đỏ đã được gắp và nâng.
\end{minipage}
\end{center}
\begin{center}
\begin{minipage}[t]{0.485\linewidth}\centering
\Anh{results/validation_T6.after_color_crop.png}
\small\textbf{(c)} T6: A--đỏ, B--vàng, C--xanh.
\end{minipage}\hfill
\begin{minipage}[t]{0.485\linewidth}\centering
\Anh{results/validation_student_P2.after_color_crop.png}
\small\textbf{(d)} T7 theo MSSV: A--vàng, B--đỏ, C--xanh; robot về home.
\end{minipage}
\end{center}
\smallskip
Ảnh màu lấy trực tiếp từ camera Gazebo ở các mốc được ghi trong JSON và chỉ
cắt vùng bàn để nhìn rõ robot, vật và vùng đích; ảnh toàn cảnh gốc vẫn nằm
trong \texttt{results/}. Ba vùng trên bàn từ phía gần đến phía xa lần lượt là
A, B, C. T6 và T7 dùng hai ánh xạ khác nhau; (d) là kết quả của \emph{một} câu lệnh
theo MSSV. Các file JSON lưu câu lệnh, phản hồi LLM, kết quả từng skill,
trạng thái MoveIt và tọa độ trước/sau.

\newpage
\section{Nhiệm vụ cá nhân theo MSSV}
Họ tên và MSSV trên bìa được lấy từ mẫu sinh viên cung cấp. Cùng MSSV
\MSSV{} đã được ghi vào cấu hình sinh viên. Hai chữ số cuối là $XX=74$, do đó
\[
P=XX\bmod 6=74\bmod 6=2.
\]
Theo quy ước của đề, trường hợp $P=2$ yêu cầu \textbf{Zone A nhận khối vàng,
Zone B nhận khối đỏ và Zone C nhận khối xanh dương}. Bảng dưới đây ghi đủ sáu
khả năng và đánh dấu hàng được chọn cho MSSV này.
\begin{center}\small
\begin{tabular}{@{}clll@{}}
\toprule $P$ & \textbf{Zone A} & \textbf{Zone B} & \textbf{Zone C}\\\midrule
0 & Đỏ & Vàng & Xanh dương\\
1 & Đỏ & Xanh dương & Vàng\\
\textbf{2 (đã chọn)} & \textbf{Vàng} & \textbf{Đỏ} & \textbf{Xanh dương}\\
3 & Vàng & Xanh dương & Đỏ\\
4 & Xanh dương & Đỏ & Vàng\\
5 & Xanh dương & Vàng & Đỏ\\\bottomrule
\end{tabular}
\end{center}

\begin{figure}[htbp]\centering
\Anh{results/code_student_assignment.png}
\caption{Hàm \texttt{assignment()} lấy đúng hai chữ số cuối MSSV, tính modulo 6 và tra ánh xạ đích.}
\label{fig:code-assignment}
\end{figure}
Hình~\ref{fig:code-assignment} cho thấy ánh xạ $P=2$ được tính từ cấu hình,
không nhập cứng kế hoạch bảy bước vào chương trình.

\textbf{Cách chạy.} Từ cảnh ban đầu, tôi nhập đúng câu lệnh
\texttt{Arrange all objects according to my student ID.} Planner nhận ánh xạ
$P=2$ trong ngữ cảnh và trả về kế hoạch bảy bước sau; validator xác nhận ba
vùng đích còn trống trước khi executor bắt đầu:
\begin{lstlisting}
pick(yellow_cube) -> place(yellow_cube, zone_a)
pick(red_cube)    -> place(red_cube, zone_b)
pick(blue_cube)   -> place(blue_cube, zone_c)
home()
\end{lstlisting}
Đây là kết quả LLM của ca T7, lưu trong
\texttt{validation\_student\_P2.json}; không phải kế hoạch nạp qua tùy chọn
\texttt{--test-plan}. Robot thực hiện đủ 7/7 skill trong
84,6~s; 19/19 chuyển động MoveIt trả mã \texttt{SUCCESS}, 12/12 đoạn Cartesian
có \texttt{fraction=1.0}, tay gắp rỗng và robot trở về home. Tọa độ tâm vật
sau tác vụ được đọc lại từ Gazebo (hệ \texttt{world}, đơn vị mét):
\begin{center}\small
\begin{tabular}{@{}lllr@{}}
\toprule\textbf{Vật / đích} & \textbf{Tâm đích} & \textbf{Tâm đo được} & \textbf{Sai số (mm)}\\\midrule
Vàng / A & $(0{,}40;-0{,}20;0{,}12)$ & $(0{,}399973;-0{,}200029;0{,}120000)$ & 0,040\\
Đỏ / B & $(0{,}40;0;0{,}12)$ & $(0{,}399982;-0{,}000025;0{,}120000)$ & 0,031\\
Xanh / C & $(0{,}40;0{,}20;0{,}12)$ & $(0{,}400006;0{,}199964;0{,}120000)$ & 0,037\\
\bottomrule
\end{tabular}
\end{center}
Các sai số trên là khoảng cách Euclid tính từ tâm đích, làm tròn đến
0,001~mm. Chúng chỉ phản ánh mô phỏng đầu hút lý tưởng. Vì ba vùng ban đầu
đều trống, ca này không cần vùng đệm để hoán đổi vật.

\section{Kết luận và liên kết bài nộp}
Hệ thống đã thực hiện thành công điều khiển một vật bằng câu lệnh tiếng Việt
và tiếng Anh, có kiểm tra kế hoạch trước khi chạy và xác nhận đích bằng trạng
thái Gazebo. Câu lệnh theo MSSV 23020774 cũng hoàn thành đủ bảy skill với
ánh xạ $P=2$; log và ảnh (d) là bằng chứng của chính lần chạy này.
Giới hạn hiện tại: đầu hút lý tưởng, chưa đánh giá trên robot thật,
chưa có vùng đệm để hoán đổi khi đích bị vật khác chiếm.

\textbf{Video demo:} tệp \path{results/ur3e_mssv_23020774_demo.mp4}
(104~s, quay trực tiếp từ camera Gazebo). Lần chạy đi kèm video lưu tại
\path{results/demo_T7_video.json}, đạt 7/7 skill, 19/19 chuyển động MoveIt
và bố trí đích $P=2$ đúng yêu cầu.\par
\textbf{Link Google Drive công khai:} \textit{[Bổ sung sau khi tải video lên.]}
\quad\textbf{GitHub:} \textit{[Bổ sung link repository Public.]}
Mã nguồn được tổ chức độc lập với bài trước; README có hướng dẫn build và chạy.

{\footnotesize\textbf{Tài liệu tham khảo:}
\href{https://github.com/UniversalRobots/Universal_Robots_ROS2_Driver}{UR ROS~2 Driver};
\href{https://github.com/ros-simulation/gazebo_ros_pkgs}{Gazebo ROS};
\href{https://github.com/decolua/9router}{9Router};
\href{https://opencode.ai/docs/en/zen/}{OpenCode Zen}.}
\end{document}
